\section{Synthetic Validation and Software Invariants}
\label{app:checks}
The synthetic scenario validates serialization and causal recorder behavior.
It is a software check, separate from the native mission comparisons.


The reproducible smoke scenario has 360 six-channel observations at five-second
intervals over three ten-minute epochs. Seed 2026 fixes Gaussian background noise.
A point spike, a sustained step, and a slow ramp are injected at declared times.
Each epoch has a 1,024-byte cap. Capture uses a first-difference threshold of 2;
the receiver uses an amplitude threshold of 4. The adaptive ring and post-trigger
window each contain three samples, with $\rho=0.2$ background reservation.
These settings are illustrative and were not selected on mission data.

\begin{table}[ht]
\centering
\footnotesize
% Generated by python/run_retention_demo.py; synthetic checks only.
\begin{tabular}{lrrrr}
\toprule
Policy & Bytes & Raw records & Raw event coverage & Receiver evidence \\
\midrule
Uniform & 3012 & 81 & 2/3 & 2/3 \\
Summaries & 2976 & 0 & 0/3 & 3/3 \\
Threshold & 276 & 5 & 2/3 & 2/3 \\
Adaptive & 1680 & 44 & 3/3 & 3/3 \\
\bottomrule
\end{tabular}

\caption{Synthetic protocol check only. Each policy has the same three-epoch
allocation of 3,072 bytes. Raw event coverage here counts events with at least
one retained raw observation, not a per-event sample fraction. Receiver evidence
for summaries is coarse bucket evidence. Actual expenditure differs by policy.}
\label{tab:smoke}
\end{table}

The artifact checks packet-byte limits, decoding, timestamp fidelity, bounded
record counts, background reservation under continuous triggering, and causal
prefix behavior. Tests also reject invalid input, malformed packets, duplicate
raw timestamps, and budget reset through reopening a finalized epoch.
Table~\ref{tab:smoke} is generated directly from the recorded experiment artifact.
It is not a statistical comparison, a diagnosis study, or a claim that the proposed
policy outperforms established mission recorders.
